\setAuthor{Jonatan Kalmus}
\setRound{piirkonnavoor}
\setYear{2019}
\setNumber{G 7}
\setDifficulty{4}
\setTopic{Varia}

\prob{Lennurada}
Lennuk vajab õhkutõusmiseks vajaliku üleslükkejõu saavutamiseks merepinna kõrgusel vähemalt $L=\SI{2}{km}$ pikkust hoorada. Kui pikka hoorada vajaks sama lennuk maailma kõrgeimal tsiviillennuväljal, mis asub merepinnast $\SI{4}{km}$ kõrgusel? Õhu tihedus merepinnal $\rho_1 = \SI{1,23}{kg/m^3}$ ning $\SI{4}{km}$ kõrgusel $\rho_2 = \SI{0,82 }{kg/m^3}$. Lennuki tiibade poolt tekitatav üleslükkejõud on võrdeline õhu tiheduse ning lennuki kiiruse ruudu korrutisega. Eeldada, et ilm on mõlemal juhul tuulevaikne ning lennuki kiirendus on kogu hoovõtu jooksul konstantne.\hint
Mõlemal juhul on lennuki vajalik üleslükkejõud sama. Selle kaudu saab siduda lennuki lõppkiiruse ning seega ka hoovõturaja pikkuse.\solu
Teame, et üleslükkejõud $F$ on võrdeline õhu tiheduse ning kiiruse ruudu korrutisega. Seega, 
$$F=C\rho v^2 \Rightarrow v^2=\frac{F}{C\rho},$$
kus $C$ on võrdetegur. Kuna lennuk alustab hoovõttu paigalseisust ühtlase kiirendusega $a$, on läbitud teepikkus $s=v^2/(2a)$. Kuna lennuki mass (ning seega ka õhkutõusuks vajaminev üleslükkejõud) ja kiirendus mõlemal juhul samad, avalduvad teepikkused vastavalt kui
$$s_1=\frac{v_1^2}{2a}=\frac{F}{2Ca\rho_1},$$ 
$$s_2=\frac{v_2^2}{2a}=\frac{F}{2Ca\rho_2}.$$   

Võrrandid omavahel jagades saame
$$\frac{s_1}{s_2}=\frac{\rho_2}{\rho_1}.$$ 
Siit 
\[
s_2 = s_1\frac{\rho_1}{\rho_2} \approx \SI{3}{km}.
\]\probend